\documentclass{article}
\usepackage{graphicx} % Required for inserting images
\usepackage{amsmath}
\usepackage{amssymb} 
\usepackage[authoryear]{natbib}
\bibliographystyle{apalike}
\usepackage[font=small,labelfont=bf]{caption}
\usepackage{authblk}
\usepackage[hidelinks]{hyperref}



\title{Notes on the MLE for Pol II speed model}
\author{Jakub Koubele}

\date{15th September, 2025}

\begin{document}	
	
	\maketitle
	
	\section{Setup and notation}
	I will use notation mostly corresponding to the BioRxiv preprint. However, I decided to explicitly index the variables to indicate the level of granularity, which led to small changes in the notation. I will be using \textbf{bold} text to denote vectors, and normal font for scalars.
	
	All parameters and observed data are for the given gene (i.e., we are fitting one model per gene). Sample annotations and library sizes are the same for the whole dataset.
	
	We have these general indices:
	
	\begin{itemize}
		\item Samples $i=1...S$.
		\item Introns $j=1...J$.
	\end{itemize}
	
	\textbf{Sample annotation}
	\begin{itemize}
		\item Vectors of explanatory variables $\mathbf{x}_i \in \mathbb{R}^p$.
		\item Library sizes $s_i$, and corresponding logarithms of them $\ell_i = \log(s_i)$.		
	\end{itemize}
	
	\textbf{Model parameters}
	
	Gene-level parameters:
	\begin{itemize}
		\item $\boldsymbol{\alpha} \in \mathbb{R}^p$, a LFC of gene expression.
		\item $c$, an intercept term for the number of exon reads.		
	\end{itemize}
	
	Intron-level parameters:
	\begin{itemize}
		\item $\boldsymbol{\beta}_j \in \mathbb{R}^p$, a LFC of Pol II speed.
		\item $\boldsymbol{\gamma}_j \in \mathbb{R}^p$, a LFC of splicing time.
		\item $d_j$, an intercept term for intron reads.
		\item $\theta_j$, a parameter related to the intercept ratio of transcription and splicing times.			
	\end{itemize}
	
	\textbf{Model predictions}
	
	Exon read counts (gene-level):	
	$$
	\mu_i =  \exp \left( c + \ell_i +\mathbf{x}_i^\mathsf{T} \boldsymbol{\alpha} \right)
	$$
	
	Intron read count (intron-level):	
	$$
	\nu_{ij} = \exp \left(d_j + \ell_i + \theta_j + \mathbf{x}_i^\mathsf{T} \left( \boldsymbol{\alpha} - \boldsymbol{\beta}_j \right) \right) + \exp \left( d_j + \ell_i + \mathbf{x}_i^\mathsf{T} \left( \boldsymbol{\alpha} + \boldsymbol{\gamma}_j \right) \right)
	$$
	
	Mixture coefficient of reads from transcribing introns (intron-level):	
	$$
	\pi_{ij} = \sigma \left( \theta_j -\mathbf{x}_i^\mathsf{T} \left( \boldsymbol{\beta}_j + \boldsymbol{\gamma}_j \right) \right)
	$$
	
	Given the mixture coefficient $\pi_{ij}$, the density of read location within an intron $j$ is predicted as:	
	$$
	p_{ij}(r) = 1 +\pi_{ij} -2\pi_{ij} r, \quad r \in [0,1]
	$$
	
	\textbf{Data}
	
	\begin{itemize}
		
		
		\item $y_{i} \in \mathbb{N}_0$ is the read count of exonic reads (gene-level).
		\item $n_{ij} \in \mathbb{N}_0$ is the read count of intronic reads (intron-level).
		\item $r_m \in [0,1]$ is the location of intronic read $m \in \{1,\ldots,n_{ij}\}$ from sample $i$ within an intron $j$ (read-level).
	\end{itemize}	
	
	\textbf{Loss function}
	
	Recall that the PMF of Poisson distribution with rate parameter $\lambda \in \mathbb{R}^+ $ is
	
	$$
	\text{Pois}(k|\lambda) = \frac{\lambda^k \exp(-\lambda)}{k!}, \quad k \in \mathbb{N}_0
	$$
	giving us the negative log-likelihood (NLL) of
	$$
	\text{NLL}_{Pois}(k|\lambda) = -\log \left( \text{Pois}(k|\lambda) \right) = \lambda - k \log(\lambda) + \text{const.}
	$$
	where the constant term $\log(k!)$ does not depend on the parameter $\lambda$.
	
	The losses for exon and intron read counts are the Poisson NLL:
	$$
	\mathcal{L}_{\text{exon}} = \sum_{i}^{S} \text{NLL}_{Pois}(y_i|\mu_i) = \sum_{i}^{S}
	\mu_i - y_i\log \left (\mu_i \right)
	$$
	$$
	\mathcal{L}_{\text{intron}} = \sum_{i}^{S} \sum_{j}^{J} \text{NLL}_{Pois}(n_{ij}|\nu_{ij}) = 
	\sum_{i}^{S} \sum_{j}^{J} \nu_{ij} - n_{ij}\log \left (\nu_{ij} \right)
	$$	
	
	The coverage loss (NLL of location density of intron reads) is:
	$$
	\mathcal{L}_{\text{coverage}} = \sum_{i}^{S}  \sum_{j}^{J} \sum_{m}^{n_{ij}} -\log \left( p(r_m|\pi_{ij})\right) = \sum_{i}^{S}  \sum_{j}^{J} \sum_{m}^{n_{ij}} -\log \left( 
	1+\pi_{ij}-2\pi_{ij}r_m \right)
	$$
	The total loss is then a sum of NLL of exon, intron and coverage components:
	$$
	\mathcal{L}_{\text{total}} = \mathcal{L}_{\text{exon}} + \mathcal{L}_{\text{intron}} + \mathcal{L}_{\text{coverage}}
	$$
	
	\section{Semi-analytical solution for special case}
	Consider a special case when $\mathbf{x}_i$ consists of a single, one-hot encoded categorical variable. We will show that in this case, the MLE solution can be stated in (almost) closed form.
	
	Without loss of generality, we will work with the case when the explanatory variable has only 2 categories. The one-hot encoded $\mathbf{x}_i$ can then be represented by a single number $x_i \in \{0,1\}$, with value $0$ for group 1 and value $1$ for group 2. Since the vectors in our notation become scalars, we will write them in a normal, not a bold font.
	
	Instead of explicitly solving for a stationary point of the loss function, we will use the
	following technique: we will independently minimize each of the components of $\mathcal{L}_{\text{total}}$, i.e.$ \mathcal{L}_{\text{exon}}$, $\mathcal{L}_{\text{intron}}$ and $\mathcal{L}_{\text{coverage}}$. We will show that a configuration of model parameters exists, which simultaneously minimizes each of these components; hence, such a solution must also minimize $\mathcal{L}_{\text{total}}$.
	
	
	\subsection{Minimizing $\mathcal{L}_{\text{exon}}$}
	
	We now define the following summary statistics for group 1 (for which we have $x_i=0$):	
	$$
	\tilde{\mu}_1 = \frac{\sum_{i: x_i=0} y_i}{\sum_{i: x_i=0} s_i}  
	$$
	In the same way, for group 2, we define:
	$$
	\tilde{\mu}_2 = \frac{\sum_{i: x_i=1} y_i}{\sum_{i: x_i=1} s_i}  
	$$
	Regarding the used notation: for the group level variables $\tilde{\mu}_1$ and $\tilde{\mu}_2$, distinguished by the tilde symbol from the sample level prediction $\mu_i$, the lower index denotes the group (1 or 2), not the sample. We will use the same notation for all group-level variables (distinguished by the tilde symbol).
	
	The values of parameters $c$ and $\alpha$ (which are the only parameters influencing $\mu_i$) that minimize $\mathcal{L}_{\text{exon}}$ can be written in terms of the group level statistics:
	$$
	c = \log \left( \tilde{\mu}_1 \right)
	$$
	
	$$
	\alpha = \log \left( \frac{\tilde{\mu}_2}{\tilde{\mu}_1}  \right)
	$$
	
	I omitted the exact derivation here, which is a rather standard MLE solution for Poisson regression with a single binary regressor and an offset (library size).	
	
	
	\subsection{Minimizing $\mathcal{L}_{\text{coverage}}$}
	We will discuss the optimal values of mixing coefficients for intron $j$, which we define as:
	$$
	\tilde{\pi}_{1j} = \arg \min \sum_{i: x_i=0}  \sum_{m}^{n_{ij}} -\log \left( 
	1+\pi_{ij}-2\pi_{ij}r_m \right)
	$$
	$$
	\tilde{\pi}_{2j} = \arg \min \sum_{i: x_i=1}  \sum_{m}^{n_{ij}} -\log \left( 
	1+\pi_{ij}-2\pi_{ij}r_m \right)
	$$
	The values of $\tilde{\pi}_{1j}$ and $\tilde{\pi}_{2j}$  do not generally have a closed-form solution. However, the loss function being minimized is convex in $\pi_{ij}$ (note that $\pi{ij}$ is the same for all samples within the group). Therefore, the values $\tilde{\pi}_{1j} \in [0, 1]$ and $\tilde{\pi}_{2j} \in [0, 1]$ can be easily determined numerically.
	Also note that $\pi_{ij}$ is not a parameter of the model, but rather an intermediate calculation.
	
	Recall that the inverse function of a sigmoid function is
	$$
	\sigma^{-1}(z) = \log \left( \frac{z}{1-z}\right)
	$$	
	This gives us the value of $\theta_j$ that leads to model prediction of $\tilde{\pi}_{1j}$ for samples in group 1:
	$$
	\theta_j = \log \left( \frac{\tilde{\pi}_{1j}}{1-\tilde{\pi}_{1j}}\right)
	$$
	The equation above is not defined for $\tilde{\pi}_{1j} \in \{ 0,1\}$. For now, we will assume that $\tilde{\pi}_{1j} \in (0,1)$, and we will return to the undefined cases later on.
	For group 2, we also want to find a model parameters that would yield the prediction of $\tilde{\pi}_{2j}$. For sample $i$ in group 2 (i.e. $x_i=1$) and intron $j$, we have:
	$$	
	\pi_{ij} = \sigma \left( \theta_j - \beta_j - \gamma_j \right) 
	$$
	Applying a sigmoid inverse function to both sides gives
	$$
	\log \left( \frac{\tilde{\pi}_{2j}}{1-\tilde{\pi}_{2j}}\right) = \theta_j - \beta_j - \gamma_j
	$$	
	Similarly above, we will assume that $\tilde{\pi}_{2j} \in (0,1)$ for now.
	
	By plugging in the optimal value of $\theta_j$ derived above, we get:
	$$
	\beta_j + \gamma_j = \log \left( \frac{\tilde{\pi}_{1j} \left( 1-\tilde{\pi}_{2j}\right)}{\tilde{\pi}_{2j} \left( 1-\tilde{\pi}_{1j}\right)}\right)
	$$
	The set of parameters that minimize $\mathcal{L}_{\text{coverage}}$ is therefore not a single point, but a line in the parameter space on which $\beta_j+\gamma_j$ fulfill the equation above (and the $\theta_j$ equals its derived optimal value).
	\subsection{Minimizing $\mathcal{L}_{\text{intron}}$}
	For sample $i$ in group 1 (so $x_i=0$) and intron $j$, we have
	$$
	\nu_{ij} = \exp(d_j + \ell_i + \theta_j) + \exp(d_j) = \exp(d_j) \exp(\ell_i) \left( \exp(\theta_j)+1 \right)
	$$
	Plugging in the optimal value of $\theta_j$ derived above, and the definition of  $(\ell_i)=\log(s_i)$, we obtain:
	$$
	\nu_{ij} = \frac{\exp(d_j) s_i}{1-\tilde{\pi}_{1j}}
	$$
	Minimization of Poisson loss for samples in group 1 can then be stated in terms of a group-level statistic. We define:
	$$
	\tilde{\nu}_{1j} = \frac{\sum_{i: x_i=0} n_{ij}}{\sum_{i: x_i=0} s_i}  
	$$
	The MLE solution for $d_j$ can then be written as
	$$
	d_j = \log \left( \tilde{\nu}_{1j} \left( 1 - \tilde{\pi}_{1j} \right) \right)
	$$
	The equation for $d_j$ follows from the MLE of the Poisson distribution with an offset; I again omitted the details for brevity.
	
	We also define the same group-level statistic for samples in group 2:
	$$
	\tilde{\nu}_{2j} = \frac{\sum_{i: x_i=1} n_{ij}}{\sum_{i: x_i=1} s_i}  
	$$
	Which gives us (by again solving for MLE of Poisson with offset):
	$$
	\tilde{\nu}_{2j} = \exp(d_j + \theta_j +\alpha -\beta_j) + \exp(d_j+\alpha+\gamma_j) 
	$$
	By plugging in the equation
	$$
	\beta_j = \log \left( \frac{\tilde{\pi}_{1j} \left( 1-\tilde{\pi}_{2j}\right)}{\tilde{\pi}_{2j} \left( 1-\tilde{\pi}_{1j}\right)}\right) - \gamma_j
	$$
	which was derived above, as well as equations for $\alpha$, $\theta_j$ and $d_j$, we get:
	$$
	\beta_j = \log \left(\frac{\tilde{\mu}_2}{\tilde{\mu}_1}\right) - \log \left(\frac{\tilde{\nu}_{2j}}{\tilde{\nu}_{1j}}\right) + \log \left(\frac{\tilde{\pi}_{1j}}{\tilde{\pi}_{2j}}\right)
	$$
	If we plug this result into the equation for $\beta_j+\gamma_j$ derived above, we also obtain
	$$
	\gamma_j = -\log \left(\frac{\tilde{\mu}_2}{\tilde{\mu}_1}\right) + \log \left(\frac{\tilde{\nu}_{2j}}{\tilde{\nu}_{1j}}\right) + \log \left(\frac{1-\tilde{\pi}_{2j}}{1-\tilde{\pi}_{1j}}\right)
	$$
	\subsection{Summary of solution for special case}
	Putting the above equations together, we obtained a solution for all model parameters written in terms of group-level statistics:
	$$
	c = \log \left( \tilde{\mu}_1 \right)
	$$	
	$$
	\alpha = \log \left( \frac{\tilde{\mu}_2}{\tilde{\mu}_1}  \right)
	$$
	$$
	\theta_j = \log \left( \frac{\tilde{\pi}_{1j}}{1-\tilde{\pi}_{1j}}\right)
	$$
	$$
	d_j = \log \left( \tilde{\nu}_{1j} \left( 1 - \tilde{\pi}_{1j} \right) \right)
	$$
	$$
	\beta_j = \log \left(\frac{\tilde{\mu}_2}{\tilde{\mu}_1}\right) - \log \left(\frac{\tilde{\nu}_{2j}}{\tilde{\nu}_{1j}}\right) + \log \left(\frac{\tilde{\pi}_{1j}}{\tilde{\pi}_{2j}}\right)
	$$
	$$
	\gamma_j = -\log \left(\frac{\tilde{\mu}_2}{\tilde{\mu}_1}\right) + \log \left(\frac{\tilde{\nu}_{2j}}{\tilde{\nu}_{1j}}\right) + \log \left(\frac{1-\tilde{\pi}_{2j}}{1-\tilde{\pi}_{1j}}\right)
	$$
	
	If any of the group-level statistics summarizing read counts is 0, some of the equations above are not defined - in this case, we simply ignore the gene or intron.
	
	We will now discuss the cases in which either $\tilde{\pi}_{1j} \in \{0, 1\}$ or $\tilde{\pi}_{2j} \in \{0, 1\}$. In this case, either $\beta_j$ or $\gamma_j$ goes to  $\pm \infty$, while the remaining parameters keep a finite value. This is similar to the case of a logistic regression with a linearly separable dataset - the parameters are pushed to $\pm \infty$, while the loss function is lower-bounded, and its minimum is not achieved for finite values of parameters.
	
	
	\section{Open question}
	
	Even though the model parameters may be easily estimated in the special case of a single categorical regressor, the solution of more general cases is unclear. In my current implementation, I am using a numerical optimizer (L-BFGS algorithm) to minimize the loss function. I would like to either provide guarantees for the optimizer, or find a counter-example in which the optimizer may get stuck in a 'bad' local minima.
	
	Several related observations and notes:
	\begin{itemize}
		\item The loss function is not quasiconvex (and hence not convex). This holds even in the simple case of a single binary regressor, for which we provided the semi-analytical solution above.
		\item The semi-analytical solution for special case relied on two facts: that we are able to easily solve each of the components $\mathcal{L}_{\text{exon}}$,  $\mathcal{L}_{\text{intron}}$ and $\mathcal{L}_{\text{coverage}}$, and that these solution have a non-empty intersection. In the more general case, this intersection is likely to be empty (however, I did not verify this yet).
		\item We are also interested in solving modified versions of the optimization problem:
		\begin{itemize}
			\item For the likelihood-ratio test (LRT), we want to solve the same problem with one element of one of the LFC vectors ($\boldsymbol{\alpha}$, $\boldsymbol{\beta}_j$, or $\boldsymbol{\gamma}_j$) enforced to always be 0. 
			\item We may want to parametrize the LFC of Pol II speed or splicing time to be gene-specific instead of intron-specific, i.e., have $\boldsymbol{\beta}_j=\boldsymbol{\beta} \quad  \forall j$ and/or $\boldsymbol{\gamma}_j=\boldsymbol{\gamma} \quad  \forall j$. 
			
		\end{itemize}
	\end{itemize}
	
	
	
	
	
	
	
	
	
	
	
	
	
	
\end{document}
